% =====================================================================
%  Abstract della tesi.
%
%  Il testo e' quello del Sommario del Lavoro di Tesi
%  (file_latex_sommario/sezioni/*.tex), adattato all'inserimento qui:
%
%    - titolazioni non numerate: prima del Capitolo 1 una \section
%      numerata uscirebbe come "0.1" e finirebbe nell'indice;
%    - equazioni e tabella in numerazione araba semplice, ripristinata
%      subito dopo, cosi' il Capitolo 1 riparte da 1.1;
%    - etichette prefissate "abs-", per non collidere con i capitoli;
%    - tabella con [H] (pacchetto float): un float differito verrebbe
%      composto dopo il ripristino dei contatori e si rinumererebbe;
%    - nessuna voce in elenco figure/tabelle;
%    - bibliografia propria rimossa: le diciassette citazioni confluiscono
%      nella bibliografia unica in fondo alla tesi (le chiavi sono gia'
%      tutte in Bibliografia.bib).
%
%  ATTENZIONE: questo file e' una COPIA del sommario, non un include.
%  Overleaf riceve solo file_latex/, quindi non si poteva puntare alla
%  cartella del sommario. Una correzione al testo va fatta due volte.
% =====================================================================

% --- pagina divisoria ------------------------------------------------
\clearpage
\thispagestyle{empty}
\phantomsection
\addcontentsline{toc}{chapter}{Abstract}
\vspace*{\fill}
\begin{center}
    {\Huge\bfseries Abstract}
\end{center}
\vspace*{\fill}
\clearpage

% --- numerazione locale ----------------------------------------------
\renewcommand{\theequation}{\arabic{equation}}
\renewcommand{\thetable}{\arabic{table}}
\let\abstractACL\addcontentsline
\renewcommand{\addcontentsline}[3]{}

\section*{Introduzione}

L'algoritmo di Shor~\cite{shor1994} fattorizza un intero composto in tempo polinomiale,
$O\!\left((\log N)^{3}\right)$, contro il tempo sub-esponenziale del \emph{general number field
sieve}~\cite{lenstra1993gnfs}, il miglior algoritmo classico noto.
È questo divario a minacciare RSA~\cite{rsa1978}, ma è asintotico e i dispositivi \emph{Noisy
Intermediate-Scale Quantum}~\cite{preskill2018} non lo realizzano: il rumore distrugge la struttura
periodica su cui l'algoritmo si fonda, le dimostrazioni restano confinate a moduli
didattici~\cite{skosana2021} e RSA-2048 richiederebbe venti milioni di qubit
fisici~\cite{gidney2021}. Restano due rimedi, la mitigazione a valle e la
correzione d'errore durante il calcolo, ai quali un numero crescente di lavori applica modelli
appresi, spesso senza un termine di confronto che isoli il contributo del modello da quello della
regola decisionale.

Il lavoro riprende un progetto preesistente --- Shor in Qiskit e una prima campagna rumorosa su
$N=15$ --- e lo rigenera dopo che una verifica per truth table ne ha invalidato il moltiplicatore
modulare controllato, il blocco che realizza $\lvert y\rangle\mapsto\lvert ay\bmod N\rangle$. Chiede quindi \emph{in quale punto della pipeline
un modello appreso migliori davvero la resa di Shor in ambiente rumoroso}, e con la stessa
tecnologia e lo stesso impianto statistico ne confronta due: un classificatore dell'istogramma di
output e un decoder addestrato sulle sindromi di un codice correttore. La risposta ha due facce, negativa a valle e positiva ma solo residuale sulla
sindrome. Il contributo è metodologico: disegno appaiato con \emph{ablazione} --- un braccio di
controllo che esegue la sola regola decisionale, senza il modello ---, separazione fra metriche
di correzione d'errore e successo di Shor, artefatti riproducibili (seed, manifest e hash del
circuito).

\section*{Materiali e Metodi}

\subsection*{Formalizzazione matematica del problema}

Per $N$ composto, dispari e non potenza di un primo, e $a$ coprimo con $N$, la fattorizzazione si
riduce alla ricerca dell'\emph{ordine}
$r=\operatorname{ord}_N(a)=\min\{k\in\mathbb{N}^{+}:a^{k}\equiv1\pmod N\}$:
se $r$ è pari e $a^{r/2}\not\equiv-1\pmod N$, allora $\gcd\!\left(a^{r/2}\pm1,N\right)$ è un fattore
non banale, con probabilità almeno $1/2$ su $a$ casuale~\cite{nielsenChuang2000}. La
parte quantistica stima $r$ per \emph{phase estimation} dell'unitario
$U_a\lvert y\rangle=\lvert ay\bmod N\rangle$: con $n$ qubit di conteggio e $M=2^{n}$, dopo la
trasformata di Fourier quantistica inversa la misura si concentra sui multipli di $M/r$, con
dispersione descritta dal nucleo di Fejér. Il post-processing $\Phi$ è deterministico: frazioni
continue di $y/M$, esatte quando $M\ge N^{2}$, poi parità e massimo comun divisore. Detto
$\mathcal{U}$ l'insieme degli $y$ da cui $\Phi$ ricava un fattore non banale,
ogni esecuzione è delimitata da due costanti in forma chiusa,
\begin{equation}
  P_{\mathrm{id}}=\!\!\sum_{y\in\mathcal{U}}\!P(y)=\frac{3}{4},
  \qquad
  P_{\mathrm{unif}}=\frac{|\mathcal{U}|}{M}=\frac{63}{256}\simeq0{,}2461 ,
  \label{eq:abs-limiti}
\end{equation}
per l'istanza canonica $N=15$, $a=7$, $n=8$, dove $r=4$ divide $M$ e $P$ è supportata sui quattro
picchi $\{0,64,128,192\}$: di questi solo $y=0$ non porta informazione sull'ordine, da cui i tre
quarti. Le due costanti fissano in anticipo la scala di lettura: $P_{\mathrm{id}}$ è la resa per
shot della distribuzione ideale; $P_{\mathrm{unif}}$ è il \emph{pavimento uniforme}, la resa attesa
quando il rumore rende equiprobabili tutte le $M$ celle, non nulla perché il post-processing accetta
$63$ esiti su $256$. Non sono limiti
universali: sono i riferimenti ideale e uniforme, e toccare il pavimento segnala la perdita del
segnale utile, non l'impossibilità di fattorizzare.

\subsection*{Modello di rumore, scenari e mitigazione a valle}

Il circuito è compilato deterministicamente (ottimizzazione 2 e seed registrato) nella base
hardware $\{r_z,\sqrt{X},X,\mathrm{CX}\}$, in cui $r_z$ è una rotazione virtuale attorno a $z$, e
il rumore è modellato su quelle porte. Dopo ogni
gate $g$ si compone un canale depolarizzante di parametro $\lambda$ con il rilassamento termico,
descritto dai tempi $T_1$ (rilassamento energetico) e $T_2$ (perdita di coerenza di fase) e
applicato in proporzione alla durata di $g$: $r_z$ ha durata nulla nel modello,
$\sqrt{X}/X$ dura 50\,ns e CX 300\,ns. La lettura è una matrice di confusione simmetrica di
parametro $p_{\mathrm{ro}}$; su $S$ shot l'istogramma normalizzato $\hat h$ è la frequenza
campionaria multinomiale della distribuzione così letta. Due preset definiscono gli scenari
confrontati: \textbf{UC1} (riferimento) con $\lambda_{1q}=10^{-3}$, $\lambda_{2q}=10^{-2}$,
$T_1=100\,\mu\mathrm{s}$, $T_2=80\,\mu\mathrm{s}$, $p_{\mathrm{ro}}=0{,}02$; \textbf{UC2} (stress)
con $\lambda_{1q}=5\times10^{-3}$, $\lambda_{2q}=5\times10^{-2}$, $T_1=50\,\mu\mathrm{s}$,
$T_2=30\,\mu\mathrm{s}$, $p_{\mathrm{ro}}=0{,}05$: errori di gate cinque volte più intensi e tempi
di coerenza dimezzati. Sono scenari uniformi illustrativi, non calibrazioni di una QPU. Come
termine di paragone a valle si valuta la \emph{zero-noise extrapolation}~\cite{temme2017}: il
circuito gira con il rumore scalato di $\lambda=1,2,3$ e l'istogramma è estrapolato a
$\lambda\to0$ (lineare su due livelli, Richardson su tre), al prezzo di altrettante repliche.

\subsection*{Politiche decisionali, modello appreso e stimando statistico}

Su $\hat h$ si definisce l'ordinamento totale $\prec$ per conteggio decrescente, con pareggi risolti
da una priorità deterministica seeded (SHA-256); la politica
$\pi_K$ prova i primi $K$ candidati secondo $\prec$. Il Metodo~1 è $\pi_1$, l'ablazione è $\pi_4$, il
Metodo~2 applica un classificatore $g_\theta(\hat h)\in\{0,1\}$ ed esegue $\pi_4$ solo se il filtro
è positivo. Il classificatore riceve $\hat h$ e l'etichetta binaria
$Y_1=\mathbf{1}[\pi_1\ \text{ha successo}]$; è addestrato su $2\,000$ istogrammi per scenario,
campionati entro il $\pm50\%$ attorno al preset, con partizione stratificata $60/20/20$. Fra random
forest,
\emph{support vector machine} a kernel RBF calibrata e percettrone multistrato, la selezione per F1
sulla validation sceglie in entrambi gli scenari la \emph{support vector machine}, valutata una
sola volta sul test. Per ogni replica si osserva il numero di iterazioni
$\mathcal{M}^{(\pi)}=\min\{m\ge1:\pi\ \text{riesce all'iterazione}\ m\}$, posto uguale a
$\mathcal{M}_{\max}+1$ in caso di fallimento, con $\mathcal{M}_{\max}=50$: una sentinella che
codifica esplicitamente la censura. Il carattere calligrafico tiene il conteggio delle iterazioni
distinto dalla dimensione $M=2^{n}$ del registro.
Lo stimando dichiarato a priori, $\rho=\bar{\mathcal{M}}_1/\bar{\mathcal{M}}_2$, resta la misura di
effetto; l'inferenza è però affidata alle differenze appaiate
$D_i=\mathcal{M}_i^{(A)}-\mathcal{M}_i^{(B)}$, perché un rapporto di medie su una variabile
discreta e censurata non ammette test esatto: test dei ranghi con segno di Wilcoxon (zeri
trattati con il metodo di Pratt), $H_0$ di simmetria attorno a zero, $\alpha=0{,}05$ e correzione di
Holm entro famiglia. Poiché ogni esecuzione Monte Carlo dà un istogramma diverso anche a parametri
identici, le tre strategie ne ricevono uno identico a ogni coppia (replica, iterazione): $D_i$
confronta le regole sulla stessa realizzazione, senza variabilità spuria. Un audit separato scarta
la variante $Y_{16}$ dell'etichetta --- successo entro i primi sedici candidati --- perché
degenerata a classe unica.

\subsection*{Correzione d'errore: formalismo, metriche e decodifica appresa}

Un codice stabilizzatore protegge l'informazione misurando ripetutamente osservabili che lasciano
invariati gli stati codificati (il sottogruppo abeliano $\mathcal{S}\subset\mathcal{P}_n$ del gruppo
di Pauli): l'esito di quelle misure è la \emph{sindrome}
$\sigma(E)$ dell'errore $E$, e un decoder $\mathcal{D}$ è la regola che dalla sindrome inferisce la
correzione. Il decoder fallisce quando la correzione riporta lo stato nel codice applicandovi però
un operatore logico non banale, cioè un errore che nessuna sindrome poteva rivelare; la probabilità
di quell'evento è il \emph{logical error rate} $p_L(p,d)$, dove $d$ è la distanza del codice, il
peso minimo di un errore non rivelabile. Si simulano il codice a ripetizione, con la forma chiusa
$p_L=3p^{2}-2p^{3}$ come test di correttezza; il
codice di Steane~\cite{steane1996}, con verifica della corrispondenza sindrome$\rightarrow$errore
sui 21 casi e stima di $\gamma$ in $p_L\propto p^{\gamma}$; il surface code
ruotato~\cite{fowler2012} sotto rumore \emph{circuit-level}, decodificato per matching di peso
minimo (MWPM) e descritto dalla legge di scala
\begin{equation}
  p_L(p,d)=A\left(p/p_{\mathrm{th}}\right)^{\lfloor(d+1)/2\rfloor},
  \label{eq:abs-scaling}
\end{equation}
con $A$ e $p_{\mathrm{th}}$ adattati ai punti sotto soglia per $d=3,5,7,9$, così da stimare la
soglia indipendentemente dall'incrocio delle curve. Il
decoder appreso è un percettrone multistrato $f_\theta$ sugli stessi \emph{detection events}, le
variazioni di sindrome fra cicli consecutivi, confrontato con il matching sia come sostituto sia
nella forma ibrida
$\mathcal{D}_{\mathrm{hyb}}(\sigma)=\mathcal{D}_{\mathrm{MWPM}}(\sigma)\oplus\mathbf{1}[f_\theta(\sigma)>t]$,
in cui la rete, quando segnala che il matching sta sbagliando, ne ribalta la decisione applicando
l'operatore logico; la soglia $t$ è scelta in validazione e la significatività dal test di McNemar.

\subsection*{Proxy fenomenologico per la sensibilità di Shor}

Dopo ogni gate compilato si applica un canale di Pauli isotropo, che con probabilità $p_g$
sostituisce l'identità con un errore di Pauli non banale scelto uniformemente, ossia
$\lambda_1=\tfrac{4}{3}p_g$ e $\lambda_2=\tfrac{16}{15}p_g$ in forma depolarizzante. È indicato
$p_g$, non $p_L$: è la probabilità di un Pauli non identità \emph{per gate compilato}, non il tasso
logico per ciclo di sindrome della sezione \emph{Correzione d'errore}, e le due unità non sono convertibili senza un modello
architetturale esplicito. Su una griglia preregistrata di $p_g$ si
misura il successo per singola misura $P_s$, da cui la cumulata $P(k)=1-(1-P_s)^{k}$ su $k$
tentativi. Non è una compilazione fault-tolerant: ne risulta un'analisi di
sensibilità, non una previsione di prestazione.

\subsection*{Strumenti, verifica e validazione}

Qiskit/Aer~\cite{qiskit2023} gestisce i circuiti generali, Stim~\cite{stim2021} e
PyMatching~\cite{pymatching2021} quelli stabilizzatori e scikit-learn~\cite{scikit_learn} i modelli
appresi; ambiente e versioni sono registrati in un formato JSON versionato (input, seed, manifest e
hash della netlist). Una demo web separa esecuzione ideale e rumorosa; codice e artefatti
sono resi disponibili con la tesi (demo: \url{https://shor-demo-6knp.onrender.com}).
L'aritmetica di Beauregard~\cite{beauregard2002}, realizzata anche per $N=21$ e $N=35$, è validata
nel solo caso ideale: $N=15$ resta l'unica istanza rumorosa. La verifica è \emph{funzionale} (truth table,
azzeramento delle ancille e confronto con la phase estimation teorica: distanza di variazione
totale $1{,}38\times10^{-9}$ per $N=21$), \emph{di identità} (profondità, conteggi e hash) e
\emph{statistica} (ipotesi e griglie preregistrate, seed separati, intervalli di Wilson,
Bonferroni sulla monotonia, Holm entro famiglia e censura esplicita). L'ablazione sugli stessi dati resta il controllo causale.

\section*{Risultati}

\textbf{Fase 1.} Sul circuito verificato ($N=15$, $a=7$, profondità 412, 224 CX, hash registrato)
i classificatori raggiungono le prestazioni ipotizzate a priori (F1 $=0{,}919$ e
$0{,}923$; AUC $=0{,}965$ e $0{,}958$), ma l'ablazione appaiata mostra che il guadagno non è loro.
Su 30 repliche per scenario, con successo $30/30$ in ogni braccio, la sola ricerca $\pi_4$ porta le
iterazioni al minimo ($\bar{\mathcal{M}}=1{,}00$ contro $1{,}37$ in UC1 e $1{,}50$ in UC2;
$p_{\mathrm{Holm}}=0{,}0024$ e $0{,}0015$), mentre il Metodo~2 non migliora la baseline ($1{,}50$ e
$1{,}37$; $p_{\mathrm{Holm}}=0{,}7314$ e $0{,}0874$) e peggiora la propria ablazione
($p_{\mathrm{Holm}}=0{,}0033$ e $0{,}0096$). Nemmeno la zero-noise extrapolation migliora la
baseline, e triplica il costo di campionamento.

\textbf{Fase 2.} I codici riproducono le previsioni a priori, con la stima
di $p_{\mathrm{th}}$ dall'Eq.~\eqref{eq:abs-scaling} concorde con l'incrocio delle curve: la ripetizione a
$d=3$ accorda con $p_L=3p^{2}-2p^{3}$ entro l'errore statistico; Steane $[[7,1,3]]$
dà $\gamma=1{,}94$ e pseudo-soglia (il rumore oltre il quale, a distanza fissata, codificare non
conviene) $p\simeq0{,}08$; il surface code per $d=3$--$9$ dà
$p_{\mathrm{th}}=0{,}86\%$ ($Z$) e $0{,}81\%$ ($X$), prefattore $A=3{,}5\times10^{-2}$ e residui di
RMS $0{,}08$ in $\ln p_L$; la decodifica appresa non vince mai come sostituto del matching e, come
correttore ibrido, riduce $p_L$ del $18$--$21\%$ a $d=3$ \emph{soltanto in presenza di errori
correlati} --- crosstalk fra qubit dati adiacenti, non rappresentabile nel grafo di matching ---,
dell'$1$--$3\%$ a $d=5$ e per nulla a $d=7$.

\textbf{Chiusura quantitativa.} Nel proxy fenomenologico descritto sopra --- in unità di $p_g$, non del $p_L$
appena riportato --- gli estremi della curva di sensibilità coincidono con le costanti
dell'Eq.~\eqref{eq:abs-limiti}: il punto senza rumore ($p_g=0$) vale $0{,}7489$ [$0{,}7459$;
$0{,}7518$] e stima $P_{\mathrm{id}}=3/4$; il plateau è già raggiunto a $p_g\simeq0{,}1$ e vale
$0{,}2459$ a $p_g=0{,}5$, cioè il pavimento uniforme
$P_{\mathrm{unif}}=63/256=0{,}2461$. La
curva non descrive un algoritmo che fallisce, ma uno che degrada verso l'estrazione casuale, con
entrambi gli asintoti fissati in anticipo dalla teoria.

\section*{Discussione e Conclusioni}

La domanda di localizzazione ha una risposta a due facce. A valle dell'esecuzione l'ipotesi
principale, $\rho\gg1$, è \emph{smentita} dai suoi stessi dati: il classificatore raggiunge le
soglie di F1 e AUC fissate a priori e resta inutile, perché scarta istogrammi che $\pi_4$ avrebbe
risolto. Sulla sindrome, dove il dato conserva una struttura locale più ricca del solo istogramma
finale, la stessa tecnologia guadagna, ma solo nella forma ibrida --- la rete affianca il matching
invece di sostituirlo --- e solo finché il rumore contiene correlazioni che il decoder analitico non
rappresenta. Il criterio unificante è quindi verificabile: un modello appreso
aggiunge valore quando sfrutta informazione che il metodo di riferimento lascia inutilizzata, non
quando ne replica la regola decisionale.

L'autovalutazione riguarda il metodo: senza il braccio $\pi_4$ la riduzione da $1{,}37$ a $1{,}00$
sarebbe stata attribuita al classificatore, un falso positivo del tutto presentabile; l'audit
dell'etichetta e la verifica per truth table del moltiplicatore, che ha imposto la rigenerazione
dell'intera campagna, sono stati più informativi delle metriche di accuratezza.

I limiti delimitano ogni conclusione: il rumore della campagna su Shor è uniforme, indipendente e
stazionario, senza topologia, crosstalk, leakage o deriva; $N=15$ è un'istanza didattica con
post-processing permissivo; le stime
QEC valgono entro il modello Clifford--Pauli, indicato come potenzialmente
ottimistico~\cite{plaquette2026}; l'analisi di sensibilità è un proxy per gate, non una simulazione
fault-tolerant. Rispetto al piano, $N=21$ e $N=35$ restano validati solo idealmente ed esclusi per
costo dalle campagne rumorose, e l'accelerazione GPU prevista non è stata impiegata.

Gli sviluppi futuri discendono dai limiti --- mappatura per operazione logica, rumore
\emph{channel-first}, decoder ricorrenti o a grafo~\cite{alphaqubit2024} --- ma la ricaduta più
immediata è metodologica: ablazione appaiata, contratto interpretativo fra
unità diverse e artefatti versionati formano un protocollo riutilizzabile per valutare qualunque
mitigazione appresa del rumore quantistico.

% --- ripristino della numerazione della tesi -------------------------
\let\addcontentsline\abstractACL
\renewcommand{\theequation}{\thechapter.\arabic{equation}}
\renewcommand{\thetable}{\thechapter.\arabic{table}}
\setcounter{equation}{0}
\setcounter{table}{0}
\clearpage


